% !TeX program = pdflatex
% Build from repository root (keep intermediates outside the repository):
% latexmk -pdf -interaction=nonstopmode -halt-on-error \
%   -outdir=/tmp/opencode/phase-field-formulation docs/phase-field-formulation.tex
% Copy the resulting PDF to docs/phase-field-formulation.pdf.
\documentclass[11pt,a4paper]{article}
\usepackage[T1]{fontenc}
\usepackage{lmodern}
\usepackage[margin=24mm,headheight=14pt]{geometry}
\usepackage{amsmath,amssymb,bm,booktabs,array,longtable}
\usepackage{fancyhdr}
\usepackage[colorlinks=true,linkcolor=blue!50!black,urlcolor=blue!50!black,citecolor=blue!50!black]{hyperref}
\usepackage{xcolor}
\hypersetup{pdftitle={Current small-strain AT2 phase-field fracture formulation},pdfauthor={mfemMechanics technical documentation}}
\newcommand{\eps}{\bm{\varepsilon}}
\newcommand{\sig}{\bm{\sigma}}
\newcommand{\uu}{\bm{u}}
\newcommand{\vv}{\bm{v}}
\newcommand{\I}{\bm{I}}
\newcommand{\D}{\bm{D}}
\newcommand{\C}{\mathbb{C}}
\newcommand{\tr}{\operatorname{tr}}
\newcommand{\dev}{\operatorname{dev}}
\newcommand{\sym}{\operatorname{sym}}
\newcommand{\dd}{\,\mathrm{d}}
\newcommand{\pos}[1]{\langle #1\rangle_+}
\newcommand{\negp}[1]{\langle #1\rangle_-}
\newcommand{\code}[1]{\texttt{\small #1}}
\newcommand{\sectionpage}[1]{\clearpage\section{#1}}
\setlength{\parindent}{0pt}
\setlength{\parskip}{6pt}
\setlength{\emergencystretch}{2em}
\renewcommand{\arraystretch}{1.2}
\pagestyle{fancy}
\fancyhf{}
\fancyhead[L]{\small mfemMechanics\enspace /\enspace AT2 formulation}
\fancyhead[R]{\small Working-tree edition}
\fancyfoot[C]{\thepage}
\numberwithin{equation}{section}
\begin{document}
\hypersetup{pageanchor=false}
\begin{titlepage}
\vspace*{12mm}
{\Large\sffamily mfemMechanics\par}
\vspace{12mm}
{\huge\bfseries Current small-strain AT2\\[4pt]phase-field fracture model\par}
\vspace{7mm}
{\Large Strong and weak forms, elastic and algorithmic moduli,\\[3pt]parameters, discretization, and state evolution\par}
\vspace{12mm}
\textbf{Implementation scope.} This document describes the current phase-field
material, finite-element integrator, and staggered nonlinear solver in the
inspected working tree. Damage is denoted by $d$ throughout:
\[
 d=0\quad\text{intact},\qquad d=1\quad\text{fully damaged}.
\]
The code's \code{phi}/\code{mPhi} is exactly $d$, not $1-d$.

\textbf{Central distinction.} Mechanical stress comes from the current degraded
elastic energy. Only the phase equation replaces the current crack-driving
energy by a maximum-history field. Consequently, unloading does not generally
admit one shared potential for both implemented residuals, and monotone history
does not itself impose pointwise damage bounds or pointwise damage irreversibility.

\vfill
\begin{tabular}{@{}ll@{}}
Inspection date & 5 September 2026\\
Base Git revision & \code{528c91c09c4d0d6a95a1396eeb345a14678ce5d0}\\
Source state & Base revision \emph{plus pre-existing uncommitted changes}\\
Artifact & Editable \LaTeX{} source and compiled PDF
\end{tabular}

\vspace{6mm}
\small The revision alone does not reproduce this working-tree edition. No C++
source was changed and no numerical experiments or C++ tests were run for this
documentation task. Source correspondence, literature verification boundaries,
and checks appropriate to this artifact are stated in Section~\ref{sec:evidence}.
\end{titlepage}
\hypersetup{pageanchor=true}

\tableofcontents
\vfill
\textbf{How to use this note.} Sections 1--3 fix conventions and inputs;
Sections 4--7 derive the continuum equations and constitutive moduli;
Sections 8--10 connect them to the implemented finite-element equations and
transactional solver. Sections 11--12 give the current example specification
and an evidence/verification guide; Section 13 records reference provenance.
All equation numbers refer to this document
unless explicitly prefixed by ``Borden report.''

\sectionpage{Scope, symbols, and code correspondence}
\label{sec:symbols}
Let $\Omega\subset\mathbb{R}^{m}$ be a fixed body, with spatial dimension
$m=2$ or $3$ and outward unit normal $\bm n$. The unknowns are displacement
$\uu(\bm x)$ and scalar damage $d(\bm x)$. The formulation is quasi-static,
small-strain, isotropic brittle fracture, with no inertia, viscosity, plasticity,
temperature evolution, or crack-face contact. The loading coordinate $\tau$ is
dimensionless continuation, not physical time. Reference and current measures
coincide to the adopted small-strain accuracy:
\begin{equation}
 \eps=\sym\nabla\uu=\tfrac12(\nabla\uu+\nabla\uu^T).
\end{equation}
For $m=2$, use the three-dimensional embedding
\begin{equation}
 \eps=\begin{pmatrix}\varepsilon_{xx}&\varepsilon_{xy}&0\\
 \varepsilon_{xy}&\varepsilon_{yy}&0\\0&0&0\end{pmatrix}.
\end{equation}
This is \textbf{plane strain}, not plane stress: $\sigma_{zz}$ generally is not
zero. Traces, deviators, eigenvalues, stresses, and moduli remain 3D.

\begin{tabular}{@{}p{.20\linewidth}p{.53\linewidth}p{.20\linewidth}@{}}
\toprule Symbol & Meaning & SI units\\\midrule
$\bm x,\uu,\ell$ & Position, displacement, crack regularization length & m\\
$d,g,k$ & Damage, degradation, residual stiffness fraction & 1\\
$E,\nu$ & Young's modulus and Poisson ratio & Pa, 1\\
$\lambda,\mu,K$ & First Lam\'e modulus, shear modulus, bulk modulus & Pa\\
$\eps,t=\tr\eps$ & Infinitesimal strain and its trace & 1\\
$\varepsilon_a,\bm n_a,\bm P_a$ & Principal strain, unit eigenvector, projector $\bm n_a\otimes\bm n_a$ & 1\\
$\psi_0,\psi^\pm,H_n,H$ & Intact/split energy density; committed/current history & J/m$^3$ = Pa\\
$\sig,\sig^\pm,\C$ & Stress, split stresses, strain-to-stress modulus & Pa\\
$G_c$ & Critical energy release rate & J/m$^2$ = N/m\\
$\bm b,\bar{\bm t}$ & Body force per volume, surface traction & N/m$^3$, Pa\\
$\vv,q;\D$ & Displacement/phase test fields; symmetric strain increment & m, 1; 1\\
$\chi$ & Current history derivative branch indicator & 1\\
\bottomrule
\end{tabular}

\begin{tabular}{@{}p{.21\linewidth}p{.72\linewidth}@{}}
\toprule Notation & Implementation correspondence\\\midrule
$d$ & \code{mPhi}, \code{phi}, \code{phaseValue}, \code{setPhaseField()}\\
$G_c,\ell,k$ & \code{criticalEnergyReleaseRate}, \code{lengthScale}, \code{residualStiffness}\\
$K$ versus $k$ & \code{bulkModulus} versus \code{getK()} (returns $k$, \emph{not} $K$)\\
$H_n,H$ & \code{CommittedValue()}, \code{TrialValue()}; local \code{positiveEnergy} is overwritten with $H$ during assembly\\
$\sig,\C$ & \code{getPK2StressVector()}, \code{getRefModuli()} in small-strain mode\\
\bottomrule
\end{tabular}

\sectionpage{Elastic constants: derivation and admissibility}
\label{sec:elastic}
The undegraded isotropic elastic energy and its first derivative are
\begin{align}
 \psi_0(\eps)&=\frac{\lambda}{2}(\tr\eps)^2+\mu\eps:\eps
 =\frac K2 t^2+\mu\dev\eps:\dev\eps,\label{eq:psi0}\\
 \sig_0&=\lambda t\I+2\mu\eps,
 \qquad\dev\eps=\eps-\frac t3\I.\label{eq:sigma0}
\end{align}
Here $\bm A:\bm B=\sum_{i,j=1}^3 A_{ij}B_{ij}$. Hydrostatic strain
$\eps=a\I$ gives $\sig_0=(3\lambda+2\mu)a\I=K t\I$, hence
\begin{equation}K=\lambda+\frac{2\mu}{3}.\end{equation}
To relate $(\lambda,\mu)$ to experimentally interpretable $(E,\nu)$, consider
\emph{undegraded, three-dimensional uniaxial stress}, with free lateral
contraction: $\eps=\operatorname{diag}(a,-\nu a,-\nu a)$ and
$\sig_0=\operatorname{diag}(Ea,0,0)$. The lateral and axial equations give
\begin{align}
 0&=\lambda(1-2\nu)-2\mu\nu,\\
 E&=\lambda(1-2\nu)+2\mu=2\mu(1+\nu).
\end{align}
Therefore the constants evaluated by the material are
\begin{equation}\boxed{
 \mu=\frac{E}{2(1+\nu)},\qquad
 \lambda=\frac{E\nu}{(1+\nu)(1-2\nu)},\qquad
 K=\frac{E}{3(1-2\nu)}.}\label{eq:constants}
\end{equation}
This derivation defines the material constants; it is not the lateral constraint
used in the 2D shear example. Do not substitute plane-stress elastic constants
into the plane-strain embedding.

\textbf{Admissible elastic inputs.} The code requires finite $E>0$ and
$-1<\nu<1/2$, pointwise for coefficient-valued inputs. These conditions make
$\mu>0$ and $K>0$, so \eqref{eq:psi0} is positive definite on nonzero symmetric
strains. The limit $\nu\to1/2$ is singular and badly conditioned in a
displacement-only formulation; no mixed incompressibility treatment is present.
The limit $\nu\to-1$ also makes the conversion singular.

\textbf{Auxetic nuance.} Negative $\nu$ is allowed and implies $\lambda<0$;
positivity of the intact energy does not require $\lambda>0$. The Amor split
uses positive $K$ and $\mu$. The spectral split instead separates a
$\lambda\pos{t}^{2}/2$ term. Its split energies remain nonnegative in the
admissible elastic range, but convexity of each split, and stability of the
damaged tangent, must not be inferred from intact stability when $\lambda<0$.
For example, with mixed-sign principal strains and $t<0$, a perturbation in a
positive principal direction has spectral tangent quadratic coefficient
$\lambda+2\mu g$, which can be negative at small $g$. The code does not impose
an additional auxetic restriction; the example uses $\nu=0.3$.

For $E=210\times10^9$ Pa and $\nu=0.3$, direct substitution yields
\[
 \lambda\simeq121.153846\ \text{GPa},\qquad
 \mu\simeq80.769231\ \text{GPa},\qquad K=175\ \text{GPa}.
\]
These are calculated constants, not fitted or simulated results.

\sectionpage{Fracture inputs, degradation, and units}
\label{sec:inputs}
The material takes two elastic coefficients, a discrete split selection, and
three scalar fracture parameters. The current defaults are
\begin{center}
\begin{tabular}{@{}llll@{}}
\toprule Input & Value in shear examples & Domain & Origin in code\\\midrule
$E$ & $210\times10^9$ Pa & finite, $>0$ & example constant\\
$\nu$ & $0.3$ & finite, $(-1,1/2)$ & example constant\\
$G_c$ & $2700$ J/m$^2$ & finite, $>0$ & parameter struct\\
$\ell$ & $15\times10^{-6}$ m & finite, $>0$ & parameter struct\\
$k$ & $10^{-9}$ & finite, $[0,1)$ & parameter struct\\
split & Miehe spectral & one of three enum values & constructor/example\\
\bottomrule
\end{tabular}
\end{center}
$G_c$ describes fracture resistance, $\ell$ controls the diffuse crack width
and influences the strength response, and $k$ retains a fraction of the
degraded stiffness. None is obtained from $E$ and $\nu$ alone. They require
independent physical identification and/or a declared regularization choice;
the repository defaults are not a documented calibration to a particular
specimen. Mesh size must resolve $\ell$, but $\ell$ is not automatically equal
to a mesh size. No universal mesh-resolution ratio is established here.

The normalized quadratic degradation is
\begin{align}
 g(d)&=(1-k)(1-d)^2+k,\label{eq:g}\\
 g'(d)&=-2(1-k)(1-d),\qquad g''(d)=2(1-k).
\end{align}
Thus $g(0)=1$ exactly, $g(1)=k$, and $g'(1)=0$. This differs from the
unnormalized expression $(1-d)^2+k$. At $k=0$, fully damaged displacement
blocks may be singular; positive $k$ improves this limit but is not a guarantee
of good conditioning or a cure for rigid motions. Values outside $[0,1]$ are
accepted by the material if finite. They are not clamped, and then $g$ need not
remain between $k$ and $1$.

The AT2 crack density and corresponding fracture energy are
\begin{equation}
 \gamma_\ell(d,\nabla d)=\frac{d^2}{2\ell}
      +\frac\ell2|\nabla d|^2,\qquad
 \mathcal G_\ell[d]=\int_\Omega G_c\gamma_\ell\dd V.\label{eq:crack}
\end{equation}
The factor $1/2$ is essential. For a flat crack on an infinite normal coordinate
$z$, the ideal profile $d(z)=\exp(-|z|/\ell)$ has
$\int_{-\infty}^{\infty}\gamma_\ell\dd z=1$. The two equal terms integrate to
one unit of crack-area density together, not two.

Equations \eqref{eq:g}--\eqref{eq:crack} adapt Borden report
\cite{borden}, Section 2.2, equations (6)--(7), printed page 5, using
$c=1-d$ and $\ell=2\epsilon_{\mathrm B}$; $\epsilon_{\mathrm B}$ is the
report's length parameter, not the strain tensor $\eps$.

\textbf{Dimensional checks.} $G_c/\ell$ and $H$ both have units Pa;
$G_c\ell|\nabla d|^2$ has units J/m$^3$. A 3D mechanical residual has units N
and a phase residual has units J (variation with respect to dimensionless
damage). Integrating the same densities over a 2D area gives energy per
out-of-plane thickness, J/m. Mechanical reactions are then N/m, whereas phase
residuals are J/m. Multiply a 2D reaction by a physical thickness in metres to
obtain force in N. No physical thickness multiplier is supplied by these examples.

\sectionpage{Energy, variations, weak form, and strong form}
\label{sec:forms}
At fixed dead body force $\bm b$ and traction $\bar{\bm t}$, first consider the
\emph{instantaneous energetic} functional before the history substitution:
\begin{equation}
 \Pi[\uu,d]=\int_\Omega\left[g(d)\psi^+(\eps)+\psi^-(\eps)
                  +G_c\gamma_\ell\right]\dd V
 -\int_\Omega\bm b\cdot\uu\dd V
 -\int_{\Gamma_t}\bar{\bm t}\cdot\uu\dd S.\label{eq:potential}
\end{equation}
Each split in Sections 5--7 satisfies $\psi^++\psi^-=\psi_0$.
Prescribe $\uu=\bar{\uu}$ on $\Gamma_u$ and, if desired, $d=\bar d$ on
$\Gamma_d$. The complementary phase boundary is $\Gamma_q$.
Seek $\uu\in[H^1(\Omega)]^m$ and $d\in H^1(\Omega)$ with those traces;
tests $\vv$ and $q$ vanish on their respective essential boundaries. Mechanical
boundary partitions may be applied componentwise.

At fixed $d$, $\delta\eps=\sym\nabla\vv$, so the displacement variation is
\begin{equation}
 R_u(\vv)=\int_\Omega\sig:\sym\nabla\vv\dd V
 -\int_\Omega\bm b\cdot\vv\dd V
 -\int_{\Gamma_t}\bar{\bm t}\cdot\vv\dd S=0,
 \quad\sig=g\sig^++\sig^-,\quad\sig^\pm=\partial_{\eps}\psi^\pm.
 \label{eq:weak-u}
\end{equation}
At fixed $\uu$, differentiating each phase term in \eqref{eq:potential} gives
$g'\psi^+q$, $G_c d q/\ell$, and $G_c\ell\nabla d\cdot\nabla q$.
The \textbf{implemented history formulation} replaces $\psi^+$ by
$H=\max(H_n,\psi^+(\eps))$ in this phase variation only:
\begin{equation}
 R_d(q)=\int_\Omega\left[g'Hq+\frac{G_c}{\ell}d q
                   +G_c\ell\nabla d\cdot\nabla q\right]\dd V=0.
 \label{eq:weak-d}
\end{equation}
There is no applied phase flux in the examples. The same $\sig$ as in
\eqref{eq:weak-u} is used, even during unloading; it does not differentiate a
history-substituted mechanical energy.

Since stress is symmetric, $\sig:\sym\nabla\vv=\sig:\nabla\vv$.
Integration by parts gives
\begin{align}
 \int_\Omega\sig:\nabla\vv\dd V
 &=-\int_\Omega(\nabla\cdot\sig)\cdot\vv\dd V
      +\int_{\partial\Omega}(\sig\bm n)\cdot\vv\dd S,\\
 \int_\Omega G_c\ell\nabla d\cdot\nabla q\dd V
 &=-\int_\Omega\nabla\cdot(G_c\ell\nabla d)q\dd V
      +\int_{\partial\Omega}G_c\ell\nabla d\cdot\bm n\,q\dd S.
\end{align}
Thus, for the current constant $G_c,\ell$, the interior strong equations are
\begin{equation}\boxed{\begin{aligned}
 \nabla\cdot\sig+\bm b&=\bm0,\\
 -G_c\ell\Delta d+\left[\frac{G_c}{\ell}+2(1-k)H\right]d
                -2(1-k)H&=0.
\end{aligned}}\label{eq:strong}\end{equation}
The natural conditions are $\sig\bm n=\bar{\bm t}$ on $\Gamma_t$ and
$G_c\ell\nabla d\cdot\bm n=0$ on $\Gamma_q$; the latter is equivalent to
$\nabla d\cdot\bm n=0$. If fracture coefficients were spatially variable, the
divergence form, not a constant-coefficient Laplacian, would be required.
Body force and traction are the standard load extension derived here; the
current domain integrator assembles internal terms only and shear examples
load through prescribed displacement.

\sectionpage{Isotropic and Amor splits: stresses and moduli}
\label{sec:amor}
For a symmetric increment $\D$, define the fixed-damage modulus by
$\C[\D]=\partial_{\eps}\sig[\D]$. Let
\begin{equation}
 \pos{x}=\max(x,0),\quad\negp{x}=\min(x,0),\qquad
 h(x)=\begin{cases}1&x>0,\\1/2&x=0,\\0&x<0.\end{cases}\label{eq:h}
\end{equation}
The slopes of the positive and negative scalar parts are chosen as $h$ and
$1-h$, respectively, at the switch. Their sum is always one.

\subsection{Isotropic degradation}
All elastic energy is crack-driving:
\begin{align}
 \psi^+&=\psi_0, & \psi^-&=0,\\
 \sig^+&=\lambda t\I+2\mu\eps, & \sig^-&=\bm0,\\
 \C^+[\D]&=\lambda\tr(\D)\I+2\mu\D,&\C^-[\D]&=\bm0,\\
 \C[\D]&=g\left[\lambda\tr(\D)\I+2\mu\D\right].
\end{align}
These expressions follow from $\delta(t^2)=2t\tr\D$ and
$\delta(\eps:\eps)=2\eps:\D$. This option degrades compression as well as
tension and shear. The superscript $+$ means ``degraded/crack-driving part,''
not necessarily tensile loading.

\subsection{Amor volumetric/deviatoric split}
The implemented split, attributed contextually to Amor et al.\ \cite{amor}, is
\begin{align}
 \psi^+&=\frac K2\pos{t}^2+\mu\dev\eps:\dev\eps,
 &\psi^-&=\frac K2\negp{t}^2,\\
 \sig^+&=K\pos{t}\I+2\mu\dev\eps,
 &\sig^-&=K\negp{t}\I.\label{eq:amor-stress}
\end{align}
To differentiate the deviatoric term, use
$\delta\dev\eps=\dev\D$ and
$\dev\eps:\dev\D=\dev\eps:\D$. The split moduli are then
\begin{align}
 \C^+[\D]&=K h(t)\tr(\D)\I+2\mu\dev\D,\\
 \C^-[\D]&=K[1-h(t)]\tr(\D)\I.
\end{align}
Combining them at fixed $d$ gives the analytic tangent used by the material:
\begin{equation}\boxed{
 \C[\D]=K\{g h(t)+1-h(t)\}\tr(\D)\I+2\mu g\dev\D.}
 \label{eq:amor-tangent}
\end{equation}
At $t=0$, the volumetric coefficient is $K(g+1)/2$. In particular, at $d=0$
it is $K$, including at zero strain and pure shear. This is the current
\textbf{corrected complementary-slope branch}; assigning zero to both scalar
branch derivatives would incorrectly remove intact bulk stiffness. At damaged
zero-trace states this is a centered generalized choice, not a classical
derivative at the nonsmooth switch.

Amor preserves purely compressive volumetric energy but degrades all deviatoric
energy, even in a state dominated by compression. A volumetric/deviatoric split
is not a face-gap constraint or frictional contact law. Both split stresses and
all tangent expressions above are obtained here by differentiating the energies
implemented in \code{EvaluateResponse}; no external equation numbers are claimed
for these derivative formulas.

\sectionpage{Miehe spectral split and its analytic derivative}
\label{sec:spectral}
Use the eigendecomposition of the physical symmetric strain tensor,
\begin{equation}
 \eps=\sum_{a=1}^{3}\varepsilon_a\bm P_a
 =\bm Q\operatorname{diag}(\varepsilon_a)\bm Q^T,
 \qquad\bm P_a=\bm n_a\otimes\bm n_a,\quad\bm Q^T\bm Q=\I.
\end{equation}
Define $\eps_\pm=\sum_a\langle\varepsilon_a\rangle_\pm\bm P_a$.
The implemented split, attributed contextually to Miehe et al.\ \cite{miehe}, is
\begin{align}
 \psi^\pm&=\frac\lambda2\langle t\rangle_\pm^2
                  +\mu\eps_\pm:\eps_\pm,\\
 \sig^\pm&=\lambda\langle t\rangle_\pm\I+2\mu\eps_\pm,
 \qquad\sig=g\sig^++\sig^-.\label{eq:spectral-stress}
\end{align}
In the first variation, the eigenvalue identity
$\delta\varepsilon_a=\bm P_a:\D$ yields
$\delta\psi^\pm=\sig^\pm:\D$. The energy is differentiable through a zero
principal strain because the squared scalar parts have zero first derivative
there. Its stress derivative needs a generalized choice at that switch.

\subsection{Matrix-function derivative, not eigenvector autodiff}
Let $f(x)=\pos{x}$ and $\widehat{\D}=\bm Q^T\D\bm Q$. At simple
eigenvalues, differentiating the spectral representation gives
\begin{equation}
 \mathcal P_+[\D]:=\delta\eps_+[\D]
   =\bm Q\left(\bm L\odot\widehat{\D}\right)\bm Q^T,
 \qquad\delta\eps_-[\D]=\D-\mathcal P_+[\D].\label{eq:frechet}
\end{equation}
$\odot$ denotes entrywise multiplication. Off-diagonal entries collect the
change of eigenvectors and have slope
$(f(\varepsilon_a)-f(\varepsilon_b))/(\varepsilon_a-\varepsilon_b)$;
diagonal entries use $f'(\varepsilon_a)$. The code evaluates the continuous
same-sign limits directly:
\begin{equation}
 L_{ab}=\begin{cases}
 1/2,&\varepsilon_a=\varepsilon_b=0,\\
 1,&\varepsilon_a,\varepsilon_b\ge0\text{ and not both zero},\\
 0,&\varepsilon_a,\varepsilon_b\le0\text{ and not both zero},\\
 \displaystyle\frac{\pos{\varepsilon_a}-\pos{\varepsilon_b}}
 {\varepsilon_a-\varepsilon_b},&\varepsilon_a\varepsilon_b<0.
 \end{cases}\label{eq:slopes}
\end{equation}
These formulas also cover repeated nonzero eigenvalues without a small-gap
regularization. Within a repeated eigenspace the resulting map does not depend
on the arbitrary eigenvector basis. At a zero eigenspace $1/2$ is the centered
generalized slope, not a claim of classical differentiability in all directions.

The split and combined moduli are
\begin{align}
 \C^+[\D]&=\lambda h(t)\tr(\D)\I+2\mu\mathcal P_+[\D],\\
 \C^-[\D]&=\lambda[1-h(t)]\tr(\D)\I+2\mu(\D-\mathcal P_+[\D]),\\
 \C[\D]&=\lambda\{g h(t)+1-h(t)\}\tr(\D)\I\notag\\
 &\quad+2\mu\bm Q\left(\{g\bm L+(\bm1-\bm L)\}
                         \odot\widehat{\D}\right)\bm Q^T.\label{eq:spectral-C}
\end{align}
Here $\bm1$ is the all-ones matrix, not $\I$. Directly weighting positive and
negative slopes avoids subtractive cancellation of tiny residual stiffness in
pure tension. All these derivatives are analytic in the current material;
no autodiff of the nonunique eigenvectors is used.

\sectionpage{Voigt convention, phase coupling, and limiting checks}
\label{sec:voigt}
The repository uses engineering shear for strain and physical shear for stress:
\begin{equation}
 \bm e=\begin{bmatrix}\varepsilon_{xx}&\varepsilon_{yy}&\varepsilon_{zz}&
 2\varepsilon_{xy}&2\varepsilon_{yz}&2\varepsilon_{xz}\end{bmatrix}^T,
 \quad
 \bm s=\begin{bmatrix}\sigma_{xx}&\sigma_{yy}&\sigma_{zz}&
 \sigma_{xy}&\sigma_{yz}&\sigma_{xz}\end{bmatrix}^T.
\end{equation}
Then $\sig:\delta\eps=\bm s^T\delta\bm e$, so the matrix modulus is
$\bm C=\partial\bm s/\partial\bm e$, not the derivative with respect to a
vector containing physical shear strain. For the $J$th unit coordinate
$\bm e_J$ construct the symmetric tensor $\bm T_J$ by inverse engineering
Voigt conversion and set
\begin{equation}
 \bm C_{:J}=\operatorname{Voigt}_{\rm stress}(\C[\bm T_J]),\quad
 \bm T_4=\tfrac12(\bm e_x\otimes\bm e_y+\bm e_y\otimes\bm e_x),
 \label{eq:voigt-C}
\end{equation}
with analogous $yz$ and $xz$ shear bases. Each off-diagonal tensor entry is
$1/2$ for a unit engineering shear. Conversely, an already assembled physical
strain tensor must \emph{not} be halved again before its eigendecomposition.

At zero damage every split recovers
\begin{equation}
 \bm C_0=\begin{pmatrix}
 \lambda+2\mu&\lambda&\lambda&0&0&0\\
 \lambda&\lambda+2\mu&\lambda&0&0&0\\
 \lambda&\lambda&\lambda+2\mu&0&0&0\\
 0&0&0&\mu&0&0\\0&0&0&0&\mu&0\\0&0&0&0&0&\mu
 \end{pmatrix}.\label{eq:C0}
\end{equation}
For isotropic degradation $\bm C=g\bm C_0$. For the other splits one must
use the strain-dependent formulas, not multiply every entry of $\bm C_0$ by
$g$. The material tangents are symmetric on smooth branches because they are
energy Hessians in the work-conjugate convention; the centered switch choices
preserve symmetry. Symmetry does not imply positive definiteness in every
parameter regime.

At fixed strain the stress--damage derivative, common to all splits, is
\begin{equation}
 \partial_d\sig=g'\sig^+,
 \qquad\partial_d\bm s=g'\bm s^+,
 \qquad\partial_{\eps}\psi^+[\D]=\sig^+:\D.\label{eq:phase-stress}
\end{equation}
This is \code{getPhaseStressDerivative()}; it does not include $\chi$.
History affects the derivative of the \emph{phase} residual with respect to
displacement, not the derivative of mechanical stress with respect to damage.

\textbf{Useful limiting checks.}
\begin{itemize}
\item At $\eps=0$, all energies and stresses vanish and the intact modulus is
\eqref{eq:C0}, including its nonzero bulk and shear parts.
\item In strictly positive principal strain (and positive trace), the spectral
law gives $\sig=g\sig_0$ and $\C=g\C_0$; at $d=1$ these retain the factor $k$.
\item In strictly negative principal strain the spectral energy has $\psi^+=0$;
Amor also has $\psi^+=0$ for purely hydrostatic compression, but not for general
compressive states with nonzero deviatoric strain.
\item In a plane-strain calculation, retain $\sigma_{zz}$ in stress invariants.
The $zz,yz,xz$ rows of the strain--displacement matrix are zero; this does not
license a plane-stress condensation of the material tangent.
\end{itemize}

\sectionpage{History substitution and the four consistent Jacobian blocks}
\label{sec:jacobian}
During a trial solve, $H_n$ is fixed at each quadrature point. Define
\begin{equation}
 H(\eps;H_n)=\max(H_n,\psi^+(\eps)),\quad
 \chi=\begin{cases}1,&\psi^+(\eps)>H_n,\\0,&\psi^+(\eps)\le H_n,\end{cases}
 \quad\delta H=\chi\sig^+:\delta\eps.\label{eq:history}
\end{equation}
Equality is explicitly \textbf{inactive}. This is distinct from the centered
$1/2$ constitutive choices in \eqref{eq:h} and \eqref{eq:slopes}.
Linearizing \eqref{eq:weak-u}--\eqref{eq:weak-d} at current $(\uu,d)$ gives
\begin{align}
 \delta R_u(\vv)&=\int_\Omega\sym\nabla\vv:
       \left[\C[\sym\nabla\delta\uu]+g'\sig^+\delta d\right]\dd V,\\
 \delta R_d(q)&=\int_\Omega\left[g' q\chi\sig^+:\sym\nabla\delta\uu
       +\left(\frac{G_c}{\ell}+g''H\right)q\delta d
       +G_c\ell\nabla q\cdot\nabla\delta d\right]\dd V.
\end{align}
These are derivatives of the implemented residual, including the derivative
of the current maximum history. They are not a frozen-history approximation
on the active branch. Dead external loads have zero derivative; other load
models would require their own tangent contributions.

At quadrature point $r$, let $n_u,n_d$ be the scalar shape-function counts,
$\bm u_e\in\mathbb R^{mn_u}$, $\bm d_e\in\mathbb R^{n_d}$, and
\begin{equation}
 \bm e=\bm B\bm u_e,\qquad d=\bm N^T\bm d_e,\qquad
 \nabla d=\bm G^T\bm d_e,\qquad
 w_r=\widehat w_r\det\left(\frac{\partial\bm x}{\partial\bm\xi}\right)_r.
\end{equation}
$\bm B$ is $6\times mn_u$, $\bm N$ is $n_d\times1$, and the rows of
$\bm G\in\mathbb R^{n_d\times m}$ are physical phase shape gradients.
The element residuals and all four Jacobian blocks are
\begin{align}
 \bm R_u^e&=\sum_r w_r\bm B^T\bm s-\bm f^e_{\rm ext},\\
 \bm R_d^e&=\sum_r w_r\left[g'H\bm N
           +G_c\ell\bm G\bm G^T\bm d_e+\frac{G_c}{\ell}d\bm N\right],\\
 \bm K_{uu}^e&=\sum_r w_r\bm B^T\bm C\bm B,\label{eq:Kuu}\\
 \bm K_{ud}^e&=\sum_r w_r(\bm B^Tg'\bm s^+)\bm N^T,\\
 \bm K_{du}^e&=\sum_r w_r g'\bm N\chi(\bm s^+)^T\bm B,\\
 \bm K_{dd}^e&=\sum_r w_r\left[G_c\ell\bm G\bm G^T
                +\left(\frac{G_c}{\ell}+g''H\right)\bm N\bm N^T\right].
 \label{eq:Kdd}
\end{align}
At active points, the cross-block contributions transpose each other. At
inactive points, $\bm K_{du}=0$ while $\bm K_{ud}$ generally is not zero.
The full assembled Jacobian is therefore \textbf{generally nonsymmetric}; it
is symmetric when all relevant cross-coupling contributions are active (or
vanish). There is \textbf{no geometric stiffness}: $\bm B$ and integration
measures refer to the fixed small-strain geometry.

\sectionpage{Discretization and what irreversibility does not guarantee}
\label{sec:discrete}
\subsection{Finite-element and quadrature contracts}
Displacement is block 0 and phase is block 1. Their $H^1$ orders $p_u,p_d$ may
differ in the integrator; the supplied examples use the same order for both.
Both residual and Jacobian assemblies choose the same default rule order,
\begin{equation}p_{\rm quad}=2\max(p_u,p_d)+1.\end{equation}
An explicit \code{SetIntRule} overrides it. This is an implementation default,
not an exactness theorem for every mixed-order product, curved geometry,
coefficient variation, or spectral response. Spectral/history switches are
nonpolynomial. A convergence study must vary quadrature as well as mesh/order.
Apply the quadrature weight and transformation measure exactly once.

MFEM element displacement arrays are component-major regardless of the global
finite-element space ordering. They are mapped as an $n_u\times m$ matrix;
global-vector accesses must still respect their actual ordering. For a 2D
displacement shape $N_i^u$, the two column contributions to $\bm B$ are
\begin{equation}
 \bm B_{i,x}=\begin{bmatrix}N^u_{i,x}&0&0&N^u_{i,y}&0&0\end{bmatrix}^T,
 \quad
 \bm B_{i,y}=\begin{bmatrix}0&N^u_{i,y}&0&N^u_{i,x}&0&0\end{bmatrix}^T.
\end{equation}
The integrator constructs $\bm F=\I+\nabla\uu$ solely to call the inherited
material interface. Its material enforces small-deformation mode; this does
not introduce a finite-strain Green--Lagrange formulation.

Point storage binds the exact displacement finite-element and integration-rule
objects and caches reference gradients. Borrowed coefficient, material,
storage, and rule objects must outlive the integrator; the form owns the
integrator. Geometry, topology, partition, finite-element, or rule changes
require resetting storage. After history develops, resetting without an
admissible transfer would discard the model state. Such AMR history transfer
and restart serialization are not provided. The examples refine only before
history initialization. Mutable scratch and state are not reentrant or
thread-safe, and the present assembly is host-only, not a device kernel.

\subsection{History is not a bound-constrained damage solve}
For constant nonnegative $H$ and spatially uniform damage, the phase equation
reduces to
\begin{equation}
 d=\frac{2(1-k)H}{G_c/\ell+2(1-k)H}.\label{eq:uniform}
\end{equation}
This lies in $[0,1)$ and increases with $H$, but it is a homogeneous limiting
case, not a discrete maximum-principle proof. In particular, AT2 admits a
nonzero homogeneous response for any positive $H$ rather than an imposed
finite damage-initiation threshold.

The code guarantees nondecreasing \emph{accepted quadrature history} through
\eqref{eq:history}. It imposes neither $d_{n+1}(\bm x)\ge d_n(\bm x)$ nor
$0\le d(\bm x)\le1$ as constraints. Arbitrary FE meshes and higher-order
spaces need not preserve continuous maximum principles. The examples check
phase coefficients after acceptance and fail if bounds are violated (tolerance
$10^{-8}$ in double precision, $10^{-4}$ in single); they do not implement a
bound-constrained solve or automatic bound-violation cutback. At higher order,
coefficient bounds do not prove bounds everywhere inside elements.

Finally, replacing $\psi^+$ by $H$ in \eqref{eq:potential} would change its
mechanical derivative on unloading. It would not produce the implemented
stress. The nonsymmetric cross blocks are another direct demonstration that
the history formulation does not generally share one potential. The inherited
unavailable element-energy method must not be interpreted as zero physical
energy or used for energy-based globalization.

\sectionpage{Transactional state and the actual staggered solver}
\label{sec:solver}
\subsection{Initialize, evaluate, accept, or reject}
Each quadrature point starts with $H_0=0$ and trial history zero. The example
also starts with $\uu=0$ and $d=0$; its notch is geometric, not a prescribed
diffuse phase crack. State evolution is:
\begin{enumerate}
\item \code{BeginStep}: copy committed history $H_n$ to trial history and open
the transaction. Save the accepted unknowns before a continuation attempt.
\item Every residual or Jacobian evaluation: compute $\psi^+$ at the current
trial displacement and reset trial history to $\max(H_n,\psi^+)$. Never take
a maximum over rejected Newton iterates. Residual/Jacobian call order must not
change the constitutive result at the same fields and committed state.
\item Success: after both current residuals converge, commit the final trial
history as $H_{n+1}$ and retain the accepted unknowns.
\item Failure: rollback trial history and restore the saved unknowns. The
continuation wrapper also restores the accepted loading coordinate and
reapplies the trial boundary data after a cutback.
\end{enumerate}
Nested begins increment a depth counter. An inner successful commit does not
write committed point history until the outermost commit. A nested rejection
marks the transaction rejected and prevents the enclosing transaction from
committing; the outer completion restores history. Assembly itself never commits.

\subsection{Block Gauss--Seidel, despite the Newton name}
\code{NewtonForPhaseField} uses the two diagonal tangent blocks, not a
monolithic solve with all four blocks. One sweep is:
\begin{align}
 \bm K_{uu}\bm c_u&=\bm r_u, &\bm u&\leftarrow\bm u-\bm c_u,\\
 \text{reevaluate }(\bm r_u,\bm r_d)&\text{ at current }(\bm u,\bm d),&&\\
 \bm K_{dd}\bm c_d&=\bm r_d, &\bm d&\leftarrow\bm d-\bm c_d,\\
 \text{reevaluate }(\bm r_u,\bm r_d)&\text{ at current }(\bm u,\bm d).&&
\end{align}
The gradients are evaluated at the corresponding current fields. A block
already within tolerance can skip its update, but is not permanently marked
converged: the other update can reactivate it. A supplied algebraic RHS is
subtracted on \emph{every} residual evaluation. There is no energy line search
in this routine, despite inheritance from \code{NewtonLineSearch}.

Let $\rho_i$ be each block's first finite nonzero residual norm in the current
solve, including evaluations between field updates. It is then frozen. The
acceptance condition, checked starting at iteration zero, is
\begin{equation}
 \|\bm r_i\|\le\max(\eta_{\rm rel}\rho_i,\eta_{\rm abs}),
 \qquad i\in\{u,d\},\quad\text{simultaneously at the same state}.\label{eq:tol}
\end{equation}
An initially zero block uses only the absolute tolerance until coupling
activates it, at which point it establishes its relative reference, even if
$\eta_{\rm abs}=0$. Nonfinite residuals and exhausted iteration budgets cause
failure, not acceptance. The same numerical absolute tolerance is used for
blocks with different physical units (Section~\ref{sec:inputs}); it is an
example heuristic, not dimensional nondimensionalization.

\sectionpage{Current shear examples: complete problem inputs}
\label{sec:examples}
Both examples use the spectral split and material inputs in
Section~\ref{sec:inputs}. The mesh must be planar 2D, contain only domain
attribute 1, and have bottom/top boundary attributes 11/12. The prescribed
conditions, with $U_{\max}\ge0$, are
\begin{equation}
 \uu=\bm0\text{ on bottom},\qquad
 u_x=U_{\max}\frac{\tau}{\tau_f},\quad u_y=0\text{ on top}.
\end{equation}
All remaining mechanical boundaries are traction-free; $\bm b=0$ and
$\nabla d\cdot\bm n=0$ on the entire phase boundary. The CLI flag
\code{-disp} names a positive displacement magnitude, not a negative sign.
There is no phase essential boundary or initial diffuse crack.

\begin{center}\small
\begin{tabular}{@{}p{.36\linewidth}p{.25\linewidth}p{.28\linewidth}@{}}
\toprule Setting / flag & Serial & MPI\\\midrule
Target & \code{PhaseField\_shear} & \code{pPhaseField\_shear}\\
Mesh (under \code{data/}) & \code{crack\_square2d.msh} & \code{crack\_square2d\_quad.msh}\\
Polynomial order, \code{-o} & 1 & 1\\
Uniform refinements & 4 (\code{-r}) & 3 serial + 4 parallel (\code{-rs}, \code{-rp})\\
Local initial refinements, \code{-lr} & 0 & 0\\
$U_{\max}$, \code{-disp} & $10^{-4}$ m & $10^{-4}$ m\\
$\tau_f$, \code{-tf} & 1 & 1\\
Initial $\Delta\tau$, \code{-dt} & $2\times10^{-6}$ & $10^{-6}$\\
Maximum $\Delta\tau$, \code{-dt-max} & $10^{-4}$ & $10^{-3}$\\
Minimum $\Delta\tau$, \code{-dt-min} & $10^{-14}$ & $10^{-14}$\\
Attempt budget, \code{-steps} & 100000 & 100000\\
Maximum nonlinear sweeps & 25 & 12\\
Output interval, \code{-oi} & 20 accepted steps & 20 accepted steps\\
\bottomrule
\end{tabular}\end{center}
Fracture flags are \code{-gc}, \code{-l}, \code{-k}. Elastic $E,\nu$ and
the split are set in source, not exposed by these CLI parsers. Require finite
$0<\Delta\tau_{\min}\le\Delta\tau_0\le\Delta\tau_{\max}$, with
$\Delta\tau_{\min}<\Delta\tau_{\max}$, and finite $\tau_f>0$.
Failed solves halve the increment; successful solves grow it by at most 1.2
up to the maximum. The attempt count includes rejections, and failure to reach
$\tau_f$ is reported. Optional local refinement selectors differ between
serial and MPI; the MPI polynomial selector is a legacy geometric heuristic,
not a predicted crack path.

Both nonlinear relative tolerances are $10^{-7}$ in double precision and
$10^{-4}$ in single. Both use
\begin{equation}
 \eta_{\rm abs}=\max\!\left(r_{\min},
       100\epsilon_{\rm mach}E\max(U_{\max},10^{-12}\text{ m})\right),
\end{equation}
where $r_{\min}$ is the smallest positive normalized \code{mfem::real\_t}
value, interpreted numerically by the code. The dimensional caveat after
\eqref{eq:tol} applies. Serial uses UMFPack; MPI uses GMRES for both diagonal
blocks with relative tolerance $10^{-10}$ ($10^{-5}$ single), 2000 iterations,
and restart length 50. Default displacement BoomerAMG uses systems options;
phase BoomerAMG is scalar. \code{-uamg} selects alternative displacement AMG
modes; \code{-il 1} enables per-solve diagnostics (default 0).

\sectionpage{Evidence, parallel interpretation, and verification guide}
\label{sec:evidence}
\subsection{Serial and MPI describe the same equations, not identical defaults}
Serial uses \code{Mesh}, \code{FiniteElementSpace}, \code{GridFunction}, and
\code{BlockNonlinearForm}; MPI uses their \code{Par} counterparts and
distributed Hypre ParCSR tangent blocks. The solver unknowns are true DOFs,
not raw element/local DOFs. Each MPI rank owns its true entries;
\code{SetFromTrueDofs} reconstructs local/shared output fields. Reaction
evaluation sums unconstrained top-horizontal internal-residual true entries;
MPI reduces owned contributions once, without gathering a global solution.
Norms and other collectives run on all ranks, with CSV/progress reporting on
rank zero. A reported 2D reaction is N/m.

Serial stress output is projected into continuous $H^1$; MPI stress output is
discontinuous. Both must use the current displacement \emph{and damage} in
\code{StressCoefficient}. \code{-no-vis} disables both ParaView and reaction
CSV; \code{-od} changes the ParaView directory, not the CSV working directory.
CMake gates serial on MFEM SuiteSparse and MPI on MFEM MPI. Use the same
explicit mesh and final refinement for rank comparisons; defaults differ.

\subsection{Source evidence and existing tests (not rerun here)}
Equations were checked against the current contents of
\code{material/PhaseFieldMaterial.h/.cpp}, \code{fem/PhaseField.h/.tpp},
\code{fem/Solvers.h/.cpp}, and the serial/MPI shear sources and their CMake
file. The existing \code{docs/phase-field-fracture.md} supplies useful context,
but current code takes precedence: notably MPI now specifies 12 nonlinear
sweeps, not the 25 stated there. This PDF is a working-tree description, not a
description of the clean base commit.

The inspected \code{tests/phase\_field\_test.cpp} includes tests for engineering
shear, all-split compression behavior and tangents, the corrected Amor zero-trace
branch, intact zero-strain stiffness, repeated principal strains, tiny residual
stiffness, active/inactive four-block Jacobians, RHS handling, initially zero
residual activation, block-solver routing, and phase-aware stress output.
\code{tests/czm\_history\_test.cpp} includes deterministic trial-history and
nested lifecycle checks. The example CMake registers serial, one-rank, and
two-rank smoke tests when targets are available. Their existence is not a new
pass result: \textbf{no C++ tests, example runs, or performance measurements
were executed for this document}. Prior run counts in other documents are not
reasserted as verification of the present working tree.

For future changes, check $\bm C\bm a$ against
$[\bm s(\bm e+\delta\bm a)-\bm s(\bm e-\delta\bm a)]/(2\delta)$ and
the full Jacobian against a centered residual directional difference with
\emph{fixed committed history}. Use smooth active and inactive states;
test history equality and zero-eigenvalue generalized branches separately.
Scale absolute-plus-relative tolerances to \code{mfem::real\_t}. Element
derivative agreement does not establish mesh convergence or calibrated fracture
prediction. Full default fracture, reaction-curve agreement, mesh/order and
quadrature convergence, auxetic damaged stability, and AMR state transfer are
not established by this documentation task.

\subsection{Reproducibility record}
For numerical results record the base revision \emph{and local diff}, MFEM
version/precision/capabilities, configure preset, executable, mesh/attributes,
order/refinement, all physical inputs and units, continuation sequence,
tolerances/iteration budgets, ranks/threads/device, and output quantities.
The source \code{docs/phase-field-formulation.tex} builds the adjacent PDF
with pdf\LaTeX{} via \code{latexmk}; intermediates are kept in
\nolinkurl{/tmp/opencode/phase-field-formulation}.

\sectionpage{Reference provenance and derivation boundaries}
\label{sec:references}
\textbf{Independently verified primary source.} The Borden author report
\cite{borden} was downloaded from the University of Texas report archive and
its relevant full-text passages inspected on 5 September 2026. Specifically:
\begin{itemize}
\item Section 2.2, equations (6)--(7), printed page 5: crack functional and
normalized quadratic degradation. Adaptations here use $c=1-d$ and
$\ell=2\epsilon_{\mathrm B}$, also identified in the report's footnote 2.
\item Section 2.2, equations (10)--(13), printed page 6: current-energy stress,
history substitution in the phase equation, and natural boundary conditions.
The present mechanical equation omits the report's inertia and permits
standard dead loads derived in Section~\ref{sec:forms}.
\item Section 4.1: quasi-static shear example, plane strain and the values
$E=210$ GPa, $\nu=0.3$, $G_c=2700$ J/m$^2$. This contextual agreement is not
a calibration or a replication of its numerical results.
\end{itemize}
All of these external equation/page numbers refer to the \emph{2011 author
report}, not to the pagination or numbering of the later journal article.

\textbf{Contextual primary-article citations.} Zotero's collection resource was
discovered, but access returned connection refused. The prior local Amor
extraction mentioned in the existing Markdown was not present in
\code{/tmp/opencode} during this task. The separate Amor and Miehe primary
article full texts were therefore \emph{not independently reverified in this
pass}. Entries \cite{amor,miehe} retain the contextual attribution and
bibliographic metadata from the existing repository documentation; the Miehe
title/authors/venue are also listed in Borden's report bibliography. No exact
Amor/Miehe article equation or page is asserted here as verified evidence.

\textbf{Derivations supplied here.} The $(E,\nu)$ conversion, first variations,
integration-by-parts steps, split stress derivatives, divided-difference
construction, engineering-Voigt conversion, and four discrete Jacobian blocks
are derivations of the stated implementation equations. They are not quoted
verbatim from an external article. The repeated-eigenvalue and zero-slope rules,
inactive history equality branch, state transaction, and solver acceptance
details are implementation-specific contracts checked against code.

\begingroup\small
\begin{thebibliography}{9}
\bibitem{borden}
M. J. Borden, C. V. Verhoosel, M. A. Scott, T. J. R. Hughes, and C. M. Landis.
\emph{A phase-field description of dynamic brittle fracture}.
ICES Report 11-14, Institute for Computational Engineering and Sciences,
The University of Texas at Austin, May 2011.
\url{https://www.oden.utexas.edu/media/reports/2011/1114.pdf}.
Verified report passages are identified above. The repository lists the later
journal version in CMAME 217--220 (2012), 77--95,
\url{https://doi.org/10.1016/j.cma.2012.01.008}; that version's equations were
not independently checked here.

\bibitem{amor}
H. Amor, J.-J. Marigo, and C. Maurini.
\emph{Regularized formulation of the variational brittle fracture with
unilateral contact: Numerical experiments}.
Journal of the Mechanics and Physics of Solids 57 (2009), 1209--1229.
\url{https://doi.org/10.1016/j.jmps.2009.04.011}.
Contextual citation; verification limitation stated above.

\bibitem{miehe}
C. Miehe, M. Hofacker, and F. Welschinger.
\emph{A phase field model for rate-independent crack propagation:
Robust algorithmic implementation based on operator splits}.
Computer Methods in Applied Mechanics and Engineering 199 (2010), 2765--2778.
\url{https://doi.org/10.1016/j.cma.2010.04.011}.
Contextual citation; verification limitation stated above.
\end{thebibliography}
\endgroup
\end{document}
